\par\Needspace{6\baselineskip}\medskip
\begin{center}\hypertarget{native-171-article124}{}\pdfbookmark[2]{124.}{article-124}124.\end{center}
Wir gehen nun zur Aufzählung der geeignetsten Verfahren über, die auf den vorhergehenden Grundsätzen beruhen. Ihre Hauptzüge sind im ersten Buch dargestellt und brauchen hier nur unserem gegenwärtigen Zweck angepasst zu werden.

Das einfachste Verfahren scheint darin zu bestehen, für $x,y$ die Entfernungen des Himmelskörpers von der Erde bei zwei Beobachtungen zu nehmen, oder vielmehr entweder die Logarithmen dieser Entfernungen oder die Logarithmen der auf die Ekliptik oder den Äquator projizierten Entfernungen. Daraus erhält man nach Artikel 64, V die heliozentrischen Orte und die zugehörigen Entfernungen von der Sonne, sodann nach Artikel 110 die Lage der Bahnebene und die heliozentrischen Längen in der Bahn. Daraus sowie aus den zugehörigen Radiusvektoren und Zeiten erhält man alle übrigen Elemente mit der Lösung der in den Artikeln 85–105 ausführlich behandelten Aufgabe. Diese stellen die beiden Beobachtungen offenbar genau dar, welche Werte man auch für $x,y$ angesetzt hat. Berechnet man nun mit diesen Elementen den geozentrischen Ort zur Zeit der dritten Beobachtung, so entscheidet seine Übereinstimmung mit dem beobachteten Ort oder seine Abweichung von ihm darüber, ob die angesetzten Werte die wahren waren oder von ihnen abwichen. Der Vergleich ist zweifach. Man kann die eine Differenz, in Länge oder Rektaszension, als $X$ und die andere, in Breite oder Deklination, als $Y$ nehmen. Ergeben sich also die Werte dieser Differenzen $X,Y$ nicht von selbst $=0$, so lassen sich die wahren Werte von $x,y$ nach dem in Artikel 120 und den folgenden Artikeln beschriebenen Verfahren ermitteln. An sich ist es beliebig, von welchen beiden der drei Beobachtungen man ausgeht. Meist empfiehlt es sich jedoch, die erste und die letzte zu wählen, abgesehen von dem besonderen Fall, auf den wir gleich eingehen.

\SourceLeaf{172}Dieses Verfahren ist den meisten später zu erläuternden vorzuziehen, weil es am allgemeinsten anwendbar ist. Ausgenommen ist der Fall, in dem die beiden äußersten Beobachtungen eine heliozentrische Bewegung von 180, 360, 540 Grad und so weiter umfassen. Dann lässt sich nämlich die Lage der Bahnebene aus zwei heliozentrischen Orten nicht bestimmen (Artikel 110). Ebenso empfiehlt es sich, das Verfahren zu meiden, wenn die heliozentrische Bewegung zwischen den beiden äußersten Beobachtungen nur sehr wenig von $180^0$ oder $360^0$ und so weiter abweicht. In diesem Fall lässt sich die Lage der Bahn nicht genau bestimmen. Genauer gesagt würden schon geringste Änderungen der für die Unbekannten angesetzten Werte so große Änderungen der Bahnlage und damit auch der Werte von $X,Y$ hervorrufen, dass man diese jenen nicht mehr als proportional ansehen könnte. Es gibt hier jedoch ein Mittel zur Abhilfe. Man geht in einem solchen Fall von der ersten und der mittleren Beobachtung oder von der mittleren und der letzten aus und nimmt für $X,Y$ die Differenzen zwischen Rechnung und Beobachtung am dritten beziehungsweise ersten Ort. Läge jedoch sowohl der zweite Ort vom ersten als auch der dritte vom zweiten etwa 180 Grad entfernt, so ließe sich die Schwierigkeit auf diese Weise nicht beheben. Solche Beobachtungen verwendet man besser nicht zur Berechnung der Elemente, da sich aus ihnen der Natur der Sache nach überhaupt keine genaue Bestimmung der Bahnlage gewinnen lässt.

Außerdem empfiehlt sich dieses Verfahren dadurch, dass man leicht abschätzen kann, wie stark sich die Elemente ändern, wenn die beiden äußersten Orte festgehalten werden und der mittlere geringfügig verändert wird. So lässt sich der Grad der Genauigkeit, den man den gefundenen Elementen zuschreiben kann, einigermaßen beurteilen.
